% Written by unit TR-13.
\section{Problem formulation}
\label{sec:problem}

\paragraph{A model, its use and the validation question.} A risk model is a map $\hat\theta=f(x;\,\mathcal D_{\mathrm{dev}})$ from an input $x$ (the attributes of
an obligor, a loan, a trade, a balance sheet or a macroeconomic path) to a quantity $\theta$ that a decision consumes, estimated on development data
$\mathcal D_{\mathrm{dev}}$. The quantity is a probability of default $\PD$, a loss rate $\LGD$, an exposure $\EAD$, an expected credit loss $\ECL$, a
quantile $\VaR_q$ or a shortfall $\ES_q$ of a loss distribution, an exposure profile, a change in economic value or a survival period of a funding gap. The use
of the model is a decision $a(\hat\theta)$ with a cost: a capital charge, a provision, a limit, an approval or a refusal. The validation question is whether the
decision rests on an estimate that is technically sound, compliant with the rule that applies and appropriate to the use; Contraste states it as three lenses,
technical robustness, regulatory compliance and appropriate use in decisions (Section~\ref{sec:introduction}).

\paragraph{The tests as decisions under size and power.} A validation test turns a claim about the model into a decision on data $\mathcal D$ that the
model did not use. It has a statistic $T(\mathcal D)$, a null hypothesis $H_0$ (for a back-test of a grade, that the true default probability is the
one the model applied), an alternative $H_1$ and a rejection region $R_\alpha$ at a level $\alpha$. Its two error rates are
\begin{equation}
  \operatorname{size}(T,\alpha)=\sup_{\vartheta\in H_0}\Prob_\vartheta\bigl(T\in R_\alpha\bigr),\qquad
  \operatorname{power}(T,\alpha;\,\delta)=\Prob_{\vartheta_\delta}\bigl(T\in R_\alpha\bigr),
  \label{eq:problem-size-power}
\end{equation}
where $\vartheta$ indexes the data-generating process and $\vartheta_\delta\in H_1$ is a process whose defect has severity $\delta$. A test whose size exceeds
$\alpha$ raises false alarms that cost a validator's time and a bank's credibility; a test whose power is low against a defect that matters passes a model that
should fail. Both are properties of the test on a class of data, not of the formula alone, and neither is stated by the instruments that prescribe the tests
\citep{credit-scoring:ecb2019,credit-scoring:bcbswp14}. The product therefore treats every test as an instrument to be calibrated: it estimates the two rates of
\cref{eq:problem-size-power} on generators whose truth is known, with $B$ independent repetitions each carrying one random stream, so that
\begin{equation}
  \widehat{\operatorname{size}}=\frac{1}{B}\sum_{b=1}^{B}\indic\{T_b\in R_\alpha\},\qquad
  \operatorname{se}=\sqrt{\frac{\widehat{\operatorname{size}}\,(1-\widehat{\operatorname{size}})}{B}},
  \label{eq:problem-mc}
\end{equation}
and each rate is reported with its Monte Carlo standard error \citep{proposed:morris2019} and its Wilson interval \citep{ratings:brown2001}. The light of a test is the
only place where a convention enters, and it is a function of the p-value or of a metric (\cref{eq:vt-light}); the 5 and 1 percent levels are conventions, not
regulatory limits.

\paragraph{The estimation problems behind the tests.} The tests judge estimates, and the estimates have their own error. For a family of models with a
reference estimator $\hat\theta$ and a known truth $\theta^\ast$ on a generator, the product measures the bias $\E[\hat\theta]-\theta^\ast$ and the root mean squared
error over independent histories, because a model that is unbiased on average can still be useless on the one history a bank has. Where the truth is unknown, as on the real
outcomes of the credit-scoring cases, the product reports the discrimination and calibration of each model on a locked holdout or an out-of-time cohort, with a
confidence interval, and compares models by a paired test (the DeLong test of the area under the receiver operating characteristic curve, \citealt{credit-scoring:delong1988}).
The assumptions of each family are stated where its method is stated: the one-factor Vasicek structure of the capital and rating sections, the Markov structure of
the transition matrices, the stationarity of the loss processes, and the exchangeability of the calibration and test windows of the conformal methods.

\paragraph{What the product claims and what it does not.} The scope is set by the design document and by its ten non-goals. A passed test is evidence, not
compliance: Contraste certifies, approves or declares no model compliant with any supervisory rule, and it represents no real bank; the synthetic institution
Banco Andino and every other synthetic institution is labelled so, and the real balance-sheet aggregates used are public aggregates. The product scores
nobody and makes no production decision about a person; the HMDA work measures a disparity in public decisions. It gives no investment advice and no forecast of
a market. It hosts no private data: the bring-your-own-data door runs in the visitor's browser and does not upload. It reproduces no proprietary dataset, and the sources
that forbid it are linked and never mirrored. It claims no general superiority of machine learning over logistic regression, of deep learning over boosting on tabular credit
data, or of a neural value-at-risk model over a classical estimator, unless the held-out comparison of the case shows it. It is no real-time risk system, and its governance
module records the product's own models. It validates generative or agentic models only as a documented practice note, and it declares no affiliation with any institution named in its sources.

\paragraph{The decision the report supports.} Given a model, a case and a battery of tests with measured error rates, the report asks what a reader may conclude. A red
light on a test whose size is held and whose power against the relevant defect is high is evidence of a defect; a green light on a test with low power is
weak evidence of soundness; and a light on a test that does not hold its size says little in either direction. The findings register (Section~\ref{sec:results}) turns
the lights into findings with a severity, and the opinion level of a case follows from a stated rule on the open findings. Each statement of this kind is made with the
case, the test and the error rate that sustain it.
